\documentclass[10pt,a4paper]{article}
\usepackage[margin=1in]{geometry}
\usepackage{hyperref}
\hypersetup{
    colorlinks=true,
    linkcolor=blue,
    urlcolor=blue,
    citecolor=blue
}
\usepackage{enumitem}
\usepackage{titlesec}
\usepackage{parskip}
\usepackage{fontawesome5}  

\titleformat{\section}{\large\bfseries}{}{0em}{}[\titlerule]

\begin{document}

\begin{center}
    {\Huge \textbf{Fernando Gutiérrez Canales}}\\
    {\Large Python Developer | Backend Engineer}\\    
    León, Mexico \quad | \quad \href{https://www.linkedin.com/in/fernando-guti\%C3\%A9rrez-canales-804876343/}{LinkedIn} \quad | \quad +52 1 477 55 18 832 \quad | \quad \href{mailto:carl.cfgc@gmail.com}{carl.cfgc@gmail.com} \\
     \href{https://github.com/Fernando-Canales}{GitHub: github.com/Fernando-Canales}
\end{center}

\vspace{1em}

\section*{Professional Summary}
Python Developer with a PhD in Exact Sciences and 5+ years of experience engineering scalable backend systems, automated workflows, and robust data pipelines. Highly proficient in integrating Python with C and Bash to optimize algorithmic performance and reduce computational overhead. Strong foundation in software engineering principles, with hands-on experience utilizing version control (Git), containerization (Docker), and relational databases (SQL) to deliver efficient, production-ready code.

\section*{Professional Experience}
\textbf{Mercor} \hfill \textbf{Göttingen (Remote)} \\
\textit{AI Prompt Engineer (LLM Logic \& Testing)} \hfill \textit{Oct 2025 -- Feb 2026}
\begin{itemize}[leftmargin=1.5em]
	\item Engineered complex logical datasets and adversarial prompts to test, debug, and fine-tune Large Language Models (LLMs), systematically identifying edge cases and improving model reasoning capacity by 30\%.
\end{itemize}

\textbf{Paris Observatory \& Max Planck Institute for Solar System Research} \hfill \textbf{Paris and Göttingen} \\
\textit{Python Developer \& Researcher} \hfill \textit{Mar 2022 -- Jul 2025}
\begin{itemize}[leftmargin=1.5em]
    \item Architected and maintained scalable automated Python pipelines to process, clean, and structure massive volumes of raw data for the PLATO space mission.
    \item Integrated C and Bash modules into existing Python codebases to optimize performance bottlenecks, drastically reducing execution time and computational resource usage.
    \item Collaborated with international teams across 15+ institutions, ensuring code quality, standardizing data architectures, and utilizing version control for seamless project integration.
\end{itemize}

\textbf{ESTEC, European Space Agency (ESA)} \hfill \textbf{Noordwijk, Netherlands}\\
\textit{Python Developer (Internship)}  \hfill \textit{Jun 2023 -- Aug 2023}
\begin{itemize}[leftmargin=1.5em]
    \item Engineered and deployed automated Python scripts to extract, validate, and parse complex parameters from hardware sensors, significantly improving the reliability of operational workflows.
\end{itemize}

\section*{Featured Development Projects}
\textbf{Plato's Closet (Open-Source Python Package)} \hfill \textbf{Python, GitHub Actions, Sphinx} \\
\textit{Lead Developer \& Maintainer} \hfill \textit{\href{https://github.com/Fernando-Canales/platos_closet}{github.com/Fernando-Canales/platos\_closet}}
\begin{itemize}[leftmargin=1.5em]
    \item Engineered and published a modular, pip-installable Python package using \texttt{setuptools}, establishing robust architecture and comprehensive technical documentation generated via Sphinx.
    \item Implemented automated Continuous Integration/Continuous Deployment (CI/CD) pipelines utilizing GitHub Actions to enforce code quality, run automated tests, and streamline release workflows.
\end{itemize}

\textbf{Ophthalmology Clinic Database Architecture} \hfill \textbf{León, Mexico} \\
\textit{Lead Developer (Independent Project)} \hfill \textit{2021 -- 2022}
\begin{itemize}[leftmargin=1.5em]
    \item Designed and implemented a centralized, relational SQL database to replace fragmented physical and digital records, establishing a robust backend architecture for the clinic.
    \item Developed automated Python scripts to query the database and generate recurring operational reports, replacing manual data entry with programmatic solutions.
\end{itemize}

\section*{Technical Skills}
\textbf{Programming Languages:} Python, C, Bash, SQL. \\
\textbf{Software Engineering \& Tools:} Git, Docker, Linux/Unix, CI/CD (GitHub Actions), Object-Oriented Programming (OOP), API Integration, Scripting \& Automation. \\
\textbf{Libraries \& Frameworks:} Pytest, Pandas, NumPy, SciPy, Scikit-learn.

\section*{Education}
\textbf{Data Analytics Bootcamp Certification} \hfill \textit{TripleTen} \hfill \textit{2026}\\
\textbf{PhD in Astrophysics} \hfill \textit{Paris Observatory \& MPS Göttingen} \hfill \textit{2025}\\
\textbf{Master's Degree in Astrophysics} \hfill \textit{University of Guanajuato, Mexico} \hfill \textit{2021}\\
\textbf{Bachelor's Degree in Physics} \hfill \textit{University of Guanajuato, Mexico} \hfill \textit{2019}

\section*{Achievements \& Publications}
\begin{itemize}[leftmargin=1.5em]
    \item Co-author of \href{https://ui.adsabs.harvard.edu/search/fq...}{peer-reviewed publications} focused on complex computational analysis and automated data processing.
    \item Recipient of full academic scholarships and graduated with Honorable Mention across all three university degrees (2019, 2021, 2025).
\end{itemize}

\end{document}
